\documentclass[openright, oneside, a4paper, article, brazil]{abntex2}
\usepackage{formularios}
% --------- %
% Metadados %
% --------- %
\hypersetup{%
	colorlinks=false,
	allbordercolors=white,
	pdftitle = {Análise — Isenção de anuidade por questões de saúde},
	pdfauthor = {Conselho Regional de Psicologia de São Paulo},
	pdfsubject = {Isenção de anuidade},
	pdfkeywords = {%
		Isenção de anuidade,
		saúde,
	}%
}
\newlength{\nomeda}
	\settowidth{\nomeda}{Nome da/o psicóloga/o: }
\newlength{\nome}
	\setlength{\nome}{\textwidth-\nomeda-0.5em}
\title{Análise — Isenção de anuidade por questões de saúde}
\author{Conselho Regional de Psicologia de São Paulo}
% --------- %
% DOCUMENTO %
% --------- %
\begin{document}
	\pagestyle{crp}
	\centering
	\Form
	\section*{ANÁLISE --- ISENÇÃO DE ANUIDADE POR QUESTÕES DE SAÚDE}%
	% -------------------------- %
	% Identificação/qualificação %
	% -------------------------- %
	\setasuspacing{\DoubleSpacing}
	% Define espaçamento duplo para que as caixas de texto não fiquem sobrepostas
	\justifying
	\preenctext{14}{nome}{\nome}{Nome da/o psicóloga/o: \hfill}\\
	\preenctext{14}{crp}{8em}{Nº de inscrição no CRP~SP: } \hfill Data de inscrição: \dataabrev{14}{1}\\
	Situação da inscrição:\hspace{1em} \caixasel{14}{ativa} ativa;\hspace{1em} \caixasel{14}{cancelada} cancelada.

	\justifying
	\subsection{Solicita isenção das anuidades}

	\begin{enumerate}
		\item \preenctext{14}{anuidade1}{47em}{\hfill}
		\item \preenctext{14}{anuidade2}{47em}{\hfill}
		\item \preenctext{14}{anuidade3}{47em}{\hfill}
		\item \preenctext{14}{anuidade4}{47em}{\hfill}
		\item \preenctext{14}{anuidade5}{47em}{\hfill}
	\end{enumerate}

\subsection{Documentos recebidos}

	\begin{enumerate}
	\item \preenctext{14}{documento1}{47em}{\hfill}
	\item \preenctext{14}{documento2}{47em}{\hfill}
	\item \preenctext{14}{documento3}{47em}{\hfill}
	\item \preenctext{14}{documento4}{47em}{\hfill}
	\item \preenctext{14}{documento5}{47em}{\hfill}
\end{enumerate}

\setasuspacing{\SingleSpace}\justifying
\subsection{Parecer técnico}
\preenctext{14}{parecer1}{\textwidth}{\null}
\preenctext{14}{parecer2}{\textwidth}{\null}
\preenctext{14}{parecer3}{\textwidth}{\null}
\preenctext{14}{parecer4}{\textwidth}{\null}
\preenctext{14}{parecer5}{\textwidth}{\null}
\vspace{0.5\baselineskip}

\noindent\small Obs.: de acordo com a resolução do CFP 20/2018 no item “5.4.2- A Diretoria do Conselho Regional deverá decidir em 30 (trinta) dias, cabendo recurso ao Plenário, no prazo de 20 (vinte) dias, se indeferido o pedido. Não havendo deliberação pela Diretoria, decorridos 30 (trinta) dias do recebimento do pedido, a interrupção temporária será considerada aprovada.


\end{document}